\documentclass[../main.tex]{subfiles}

\begin{document}
\begin{theorem}
    \label{th:1.38}
    Sia $f : G_1 \rightarrow G_2$ un morfismo di gruppi. Allora:
    \[
        \text{$f$ è iniettivo} \iff \ker(f) = \{e_1\}
    \]
    \begin{note}
        Per i morfismi di monoidi vale invece:
        \[
            \text{$f$ è iniettivo} \implies \ker(f) = \{e_1\}
        \]
    \end{note}
\end{theorem}

\begin{proof}\
    \begin{itemize}
        \item ($\implies$) Sia $f$ iniettivo. Sia $x \in \ker(f)$, allora $f(x) = e_2$ (per la definizione di $\ker$) e quindi, poiché anche $f(e_1) = e_2$ (per la condizione di morfismo), si ha che $x = e_1$ per l'ipotesi di iniettività.

        \item ($\impliedby$) Sia $\ker(f) = \{e_1\}$ e siano $x,y \in G_1$ tali che $f(x) = f(y)$.

              Allora:
              \[
                  f(x)f(y^{-1}) = e_2 \implies  f(xy^{-1}) = e_2 \implies  xy^{-1} \in \ker(f) \implies xy^{-1} = e_1 \implies x=y
              \]
              Cioè $f$ è iniettivo.
    \end{itemize}
\end{proof}

\begin{example}
    \
    \begin{itemize}
        \item $G = \mathbb{Z}_4 = \{\overline{0}, \overline{1}, \overline{2}, \overline{3}\}$, \
              \begin{itemize}
                  \item $\langle \overline{0} \rangle = \{\overline{0}\}$ sottogruppo banale $\simeq \mathbb{Z}_1$
                  \item $\langle \overline{1} \rangle = \mathbb{Z}_4$
                  \item $\langle \overline{2} \rangle = \{\overline{0}, \overline{2}\} \simeq \mathbb{Z}_2$ ($2+2=0$)
                  \item $\langle \overline{3} \rangle = \mathbb{Z}_4$ ($3 , 3+3=6=2, 3+2=5=1, 3+1=4=0$)
              \end{itemize}
              I sottogruppi di $\mathbb{Z}_4$ possono averer cardinalità $1,2,4$. L'insieme
              dei sottogruppo di $\mathbb{Z}_4$ è $\{\{\overline{0}\}, \{\overline{0}, \overline{2}\},
                  \{\overline{0}, \overline{1}, \overline{2}, \overline{3}\}= \mathbb{Z}_4\}$
        \item $G = \mathbb{Z}_6 = \{\overline{0}, \overline{1}, \overline{2}, \overline{3}, \overline{4}, \overline{5}\}$, \
              \begin{itemize}
                  \item $\langle \overline{0} \rangle = {\overline{0}}$ sottogruppo banale $\simeq \mathbb{Z}_1$
                  \item $\langle \overline{1} \rangle = \mathbb{Z}_6$
                  \item $\langle \overline{2} \rangle = \{\overline{0}, \overline{2}, \overline{4}\} \simeq \mathbb{Z}_3$
                  \item $\langle \overline{3} \rangle = \{\overline{0}, \overline{3}\} \simeq \mathbb{Z}_2$
                  \item $\langle \overline{4} \rangle = \{\overline{0}, \overline{2}, \overline{4}\} \simeq \mathbb{Z}_3$
                  \item $\langle \overline{5} \rangle = \mathbb{Z}_6$
              \end{itemize}
              I sottogruppi di $\mathbb{Z}_6$ possono averer cardinalità $1,2,3,6$. L'insieme dei sottogruppo di $\mathbb{Z}_6$ è $\{\{\overline{0}\}, \{\overline{0}, \overline{2}, \overline{4}\}, \{\overline{0}, \overline{3}\}, \{\overline{0}, \overline{1}, \overline{2}, \overline{3}, \overline{4}, \overline{5}\}= \mathbb{Z}_6\}$
    \end{itemize}
\end{example}

Trattiamo ora il caso generale: consideriamo il gruppo
\[
    \mathbb{Z}_n = (\{\overline{0},\overline{1},\ldots,\overline{n-1}\},+) \qquad n \in \mathbb{N}
\]
Sia $m \in \mathbb{N}$, $m < n$.

\begin{itemize}
    \item \textbf{Caso Base}: se $m=0$, $\langle \overline{0} \rangle = \{\overline{0}\}$ perchè sommando $\overline{0}$ a se stesso ottengo solo $\overline{0}$.
    \item \textbf{Caso Generale}: se $m > 0$, vogliamo determinare il sottogruppo generato da $\overline{m}$ ossia
          \[
              \langle \overline{m} \rangle = \{\overline{m},\overline{2m},\ldots\}
          \]
          Il numero di elementi distinti di questo insieme è l'ordine del sottogruppo $\langle \overline{m} \rangle$, cioè il \textbf{minimo} intero $z$ tale che
          \[
              \overline{zm} = \overline{0}
          \]
          Questo significa che $zm$ è multiplo di $n$, ossia $zm$ è il \textbf{minimo multiplo comune} tra $m$ e $n$, quindi:
          \[
              zm = \mcm \{m,n\} \implies z = \frac{\mcm \{m,n\}}{m}
          \]
\end{itemize}
Dimostriamo ora che $z$ è il più piccolo valore per il quale $\overline{zm} = \overline{0}$.

\begin{proof}
    Ipotizziamo per assurdo che esista $1 \leq i < z$ tale che $\overline{im} = \overline{0}$.

    Allora $im$ dovrebbe essere multiplo di $n$. Ma per definizione di $z$ sappiamo che
    \[
        zm = \mcm \{m,n\}
    \]
    e dato che $i < z$ si ha che $im < zm$ e quindi $im$ non può essere multiplo di $n$ perchè $zm$ è il minimo multiplo comune tra $m$ e $n$.

    Dunque $1 \leq i < z$, $\overline{im} \neq \overline{0}$, quindi $\overline{m}$ torna a zero solo dopo $z$ iterazioni.
\end{proof}
Dunque:
\[
    \boxed{
        \abs{\langle \overline{m} \rangle} = z = \frac{\mcm \{m,n\}}{m}
    }
\]
In particolare
\[
    \langle \overline{m} \rangle = \mathbb{Z}_n \iff \abs{\langle \overline{m} \rangle} = \abs{\mathbb{Z}_n} \iff \frac{\mcm \{m,n\}}{m} = n \iff \mcm \{m,n\} = mn \iff m \text{ e } n \text{ sono coprimi}
\]

Ossia \underline{l'insieme $\{\overline{m}\}$ genera il gruppo $\mathbb{Z}_n$ sse $m$ e $n$ sono coprimi} ($\MCD \{m,n\}=1$).

\begin{definition}[Funzione di Eulero]
    La funzione definita da
    \[
        \varphi : \mathbb{N} \setminus \{0\} \rightarrow \mathbb{N} \setminus \{0\} \quad \varphi(n) := \left|\{m \in \mathbb{N} \setminus \{0\} : m < n \text{ e } \MCD \{m,n\} = 1\}\right|
    \]
    è detta \textbf{funzione di Eulero}.

    Quindi ci sono $\varphi(n)$ elementi $\overline{m}$ tali che $\langle \overline{m} \rangle = \mathbb{Z}_n$.
\end{definition}

\begin{example}\
    $\varphi(6) = 2$ perché solo $1$ e $5$ sono coprimi con $6$ (e $1$ e $5$ generano $\mathbb{Z}_6$).
\end{example}

\begin{proposition}
    L'insieme dei sottogruppi di $(\mathbb{Z} ,+)$ è
    \[
        \boxed{
            \{n \mathbb{Z} : n \in \mathbb{N} \}
        }
    \]
\end{proposition}

\begin{proof}
    Sia $H \subseteq \mathbb{Z} $ un sottogruppo non banale (Se $H$ banale allora $H = 0\mathbb{Z}$).
    Sia $k := \min (H_{>0})$ dove
    \[
        H_{>0} := \{h \in H : h > 0\}
    \]
    Sia $h \in H_{>0}, h \neq k$. Allora $h > k$ e scrivendolo secondo la divisione euclidea dove $n$ è il quoziente della divisione $\frac{h}{k}$ e $r$ il resto:
    \[
        h = nk + r \qquad 0 \leq r < k \qquad n \in \mathbb{N}, h,k,r \in \mathbb{Z},
    \]
    Se avessimo $r \neq 0$ allora $r$ sarebbe in $H_{>0}$ e $r < k$ che contraddirebbe la definizione di $k$ data inzialmente.

    L'unico modo per avere $0 \leq r < k$ senza invalidare la minimalità di $k$ è che $r \notin H_{>0}$, cioè $r = 0$.

    Dunque:
    \[
        h = nk + 0 = nk
    \]
    Quindi tutti i sottogruppi non banali in $\mathbb{Z}$ sono della forma $h = nk$ dove $n \in \mathbb{N}$, $k \in \mathbb{Z}$.
\end{proof}

\begin{definition}[Gruppo ciclico]
    Un gruppo $G$ è detto \textbf{ciclico} se esiste $g \in G$ tale che $\langle g \rangle = G$.
    \begin{note}\
        \begin{itemize}
            \item Un gruppo ciclico è anche abeliano.
            \item Se $\abs{G} = p$ con $p$ primo, allora $G$ è ciclico.
        \end{itemize}
    \end{note}
\end{definition}
\begin{example}
    \
    \begin{itemize}
        \item $\mathbb{Z} = \langle 1 \rangle$ è ciclico
        \item $\mathbb{Z}_n = \langle \overline{1} \rangle$ è ciclico
        \item $\mathbb{Z} \times \mathbb{Z} = \langle (1,0), (0,1) \rangle$ \underline{non} è ciclico,

              infatti in $\mathbb{Z} \times \mathbb{Z} $, se $(a,b) \in \mathbb{Z} \times \mathbb{Z} $
              \[
                  \langle (a,b) \rangle = \{(ka,kb) : k \in \mathbb{Z} \} = \{(x,y) : a \text{ divide } x,b \text{ divide } y\} \subsetneq \mathbb{Z} \times \mathbb{Z}
              \]
        \item $\mathbb{Z}_2 \times \mathbb{Z}_2$ \underline{non} è ciclico, infatti, in $\mathbb{Z}_2 \times \mathbb{Z}_2$ si ha:
              \begin{itemize}
                  \item $\langle (\overline{0},\overline{0}) \rangle = \{(\overline{0},\overline{0})\}$
                  \item $\langle (\overline{0}, \overline{1}) \rangle = \{\overline{0}\} \times \mathbb{Z}_2$
                  \item  $\langle (\overline{1}, \overline{0}) \rangle = \mathbb{Z}_2 \times \{\overline{0}\}$
                  \item $\langle (\overline{1}, \overline{1}) \rangle = \{(\overline{0},\overline{0}),(\overline{1},\overline{1})\}$
              \end{itemize}
              Quindi nessun elemento di $\mathbb{Z}_2 \times \mathbb{Z}_2$ genera $\mathbb{Z}_2 \times \mathbb{Z}_2$.
    \end{itemize}
\end{example}

\begin{theorem}[di isomorfismo per gruppi abeliani]
    Sia $f: G_1 \rightarrow G_2$ un morfismo di gruppi abeliani. Allora esiste un morfismo iniettivo $\varphi : \nicefrac{G_1}{\ker \varphi} \rightarrow G_2$ tale che il seguente diagramma è commutativo:
    \begin{equation*}
        \begin{tikzcd}
            G_1 \arrow[r, "f"] \arrow[d, "\pi"'] & G_2 \\
            \nicefrac{G_1}{\ker(f)} \arrow[ru, "\varphi"']
        \end{tikzcd}
    \end{equation*}
    In particolare, $\nicefrac{G_1}{\ker(f)} \simeq \Im(f)$.
    \begin{note}\
        \begin{itemize}
            \item "Commutativo" cioè $f = \varphi \circ \pi$
            \item $\varphi ([g]) = f(g)$
            \item Se $f$ è anche suriettiva allora $\nicefrac{G_1}{\ker(f)} \simeq G_2$
        \end{itemize}
    \end{note}
\end{theorem}

\begin{proof}\
    \begin{itemize}
        \item  L'assegnazione $[g] \mapsto f(g), \forall g \in  G$, definisce una funzione
              \[
                  \varphi : \nicefrac{G_1}{\ker(f)} \rightarrow G_2
              \]
              vediamo come questa funzione è ben definita: se $g' \sim g$, ossia $[g]  = [g']$ allora vogliamo $\varphi([g]) = \varphi([g'])$.

              Se $[g] = [g']$ per definizione dell'insieme quoziente significa che esiste un elemento $h \in \ker(f)$ tale che
              \[
                  g = g' + h \qquad h \in \ker(f)
              \]
              Dunque, applicando $f$ ad entrambi i membri dell'equazione otteniamo
              \[
                  f(g) = f(g' + h) = f(g') + \underbrace{f(h)}_{0_2} = f(g')
              \]
              Poichè $h \in \ker(f)$, per definizione di nucleo si ha $f(h) = 0_2$ (identità di $G_2$). Abbiamo ottenuto che $\varphi([g]) = \varphi([g'])$ quindi \textbf{la funzione è ben definita}.

        \item Se $f$ è morfismo di gruppi, anche $\varphi$ lo è, infatti per essere morfismo deve valere
              \[
                  \varphi([g] + [g']) = \varphi([g]) + \varphi([g'])
              \]
              ma poichè $f$ è morfismo di gruppi vale:
              \[
                  f(g+g') = f(g) + f(g')
              \]
              ricordando la relazione tra $f$ e $\varphi$: $\varphi([g]) = f(g)$, otteniamo
              \[
                  \varphi([g] + [g']) = f(g+g') = f(g) + f(g') = \varphi([g]) + \varphi([g'])
              \]
              quindi \textbf{$\varphi$ è un morfismo di gruppi}.
        \item Mostriamo che $\varphi$ è iniettivo sfruttando il teorema \ref{th:1.38}, per definizione del nucleo abbiamo
              \[
                  \ker(\varphi) = \{[g] \in \nicefrac{G_1}{\ker(f)} : \varphi([g]) = 0_2\}
              \]
              questo significa che $[g]$ sta nel nucleo se e solo se
              \[
                  \varphi([g]) = f(g) = 0_2
              \]
              ma questo equivale a dire che $g \in \ker(f)$. Qual è la classe di $g$ nel quoziente?
              \[
                  [g] = g + \ker (f) = \{g + h : h \in \ker(f)\}
              \]
              Ma dato che $g \in \ker(f)$ possiamo riscrivere la classe come
              \[
                  [g] = \ker(f)
              \]
              che è precisamente la classe neutra (identità del quoziente)
              \[
                  [0_1] = 0_1 + \ker(f) = \ker(f)
              \]
              Quindi tutti gli elementi di $\ker(f)$ formano \textbf{un'unica classe} nel quoziente, che è a classe neutra del gruppo $\nicefrac{G_1}{\ker(f)}$ indicata come $[0_1]$ quindi il nucleo di $\varphi$ è
              \[
                  \ker(\varphi) = \{[0_1]\}
              \]
              Essendo il nucleo di $\varphi$ banale, \textbf{la funzione è iniettiva}.

        \item $\varphi : \nicefrac{G_1}{\ker (f)} \to \Im (f)$ è anche suriettiva,
              cioè
              \[
                  \forall y \in \Im(f) \quad \exists [g] \in \nicefrac{G_1}{\ker(f)} : \varphi([g]) = y
              \]
              prendendo qualsiasi elemento $y \in \Im(f)$, per definizione di $\Im (f)$ $\exists g \in G_1 : f(g) = y$, ma allora, sempre dalla relazione tra $f$ e $\varphi$ otteniamo:
              \[
                  f(g) = \varphi([g]) = y
              \]
              quindi per ogni elemento di $\Im(f)$ esiste una classe $[g] \in \nicefrac{G_1}{\ker(f)}$ che lo mappa su di esso, cioè $\varphi$ è    \textbf{suriettivo}.
    \end{itemize}
    $\varphi$ è sia iniettivo che suriettivo, quindi è un \textbf{isomorfismo}.
    Dunque $\nicefrac{G_1}{\ker(f)} \simeq \Im(f)$.
\end{proof}
\end{document}